\section{Distributional Fidelity of the Training Signal} \label{sec:14-introduction-fidelity}
\begin{foldable}
  The restrictions described in the previous subsections do not merely reduce the quantity of available supervision, but instead reshape its statistical character.
  A student trained under a limited data constraint, or a model fine-tuned on a compressed corpus can only learn from whatever training signal is actually present.
  If that signal fails to faithfully represent the distribution the teacher was trained on, then the student inherits a systematically distorted picture of the teacher's knowledge.
  The term \textit{distributional fidelity} names this property: the degree to which the training signal available to the student covers the target distribution in a representative and unbiased way.
  When distributional fidelity is low, distillation fails because the information it needs is not present in the signal it receives.
\end{foldable}

\begin{foldable}
  In model distillation, distributional fidelity shows up as a coverage problem that worsens as data access diminishes.
  Gaps in coverage correspond directly to gaps in the student's knowledge domain: regions of the teacher's input distribution that the student never trains on remain unknown to it, regardless of model capacity.
  In the few-shot setting, a handful of real samples cannot span the full diversity of the teacher's input distribution, and synthetic data generated from them inherits the same distributional limits.
  Generative models can extend synthesis beyond these limits, but without external guidance they remain tethered to the modes the few-shot set already represents.
  Diversity of the generated set, specifically its ability to reach regions the few-shot examples do not represent, is therefore the binding challenge.
  In the data-free setting, the problem is more severe: without any real-data anchor, synthetic images serve as a substitute, seeding synthesis with structured diversity before any teacher query.
  Yet this does not guarantee coverage: the central challenge is ensuring that synthetic samples collectively capture \textit{diverse patterns and semantics} of the teacher's knowledge domain, rather than merely satisfying its hard labels.
\end{foldable}

\begin{foldable}
  The black-box constraint compounds the limited data setting by further degrading the signal quality.
  In practice, black-box access exists on a spectrum: some services return a truncated probability distribution over the top-$k$ classes, while others return only a single integer identifying the predicted class.
  Even the softer end of this spectrum degrades distributional fidelity: a truncated distribution discards the long tail of inter-class relationships that the teacher's full softmax encodes.
  At the hard-label extreme, the distilled prediction distribution collapses to a one-hot spike, carrying only the identity of the predicted class and none of the relational structure that distinguishes \ac{KD} from ordinary supervised training.
  The central challenge is therefore to recover knowledge in the form of soft labels, the inter-class relational structure that grounds the original \ac{KD} formulation~\cite{hinton2015distilling}, from a teacher interface that deliberately withholds it.
\end{foldable}

\begin{foldable}
  Dataset distillation tackles distributional fidelity from a different angle: preserving knowledge for downstream utility.
  Compared with prototype and coreset methods, dataset distillation approximates the original data with synthetic samples, free of the constraint of the existing data support.
  In principle, this can raise the fidelity and downstream quality of the distilled data.
  At extreme compression ratios, however, preserving fidelity is still hard: each distilled sample must carry a disproportionately higher informational load.
  A principled method must therefore perform \textit{knowledge estimation} on each sample and ensure that the distilled set collectively covers the full distribution, rather than over-representing some regions and under-representing or omitting others.
  The optimization objective must therefore be \textit{global}: selecting samples independently by any per-sample score cannot guarantee coverage, since it ignores relationships between chosen samples and leaves redundant or unrepresented regions uncorrected.
  The \textit{discrete} nature of text also prevents the original gradient-based optimization from transferring directly.
  This constraint demands a discrete optimization method that handles text tokens without sacrificing human-readability.
\end{foldable}

\begin{foldable}
  All three settings exhibit the same underlying problem: failure modes that undermine distributional fidelity.
  In model distillation, coverage deficit arises through generative collapse in \textit{black-box few-shot} or lack of knowledge anchor \textit{black-box data-free}.
  In dataset distillation, it arises through compression and interpretability requirements of \textit{text datasets}.
  Across all cases, the training signal fails to faithfully represent the target distribution: it over-represents some regions and under-represents or entirely omits others.
  The diagnostic question is identical in every setting: does the data the student actually trains on give sufficient, representative coverage of the knowledge the teacher possesses?
  Fidelity deficit is the root cause of these failures, not a symptom.
  No algorithmic refinement can recover knowledge that is absent from the training signal.
  Recognising this unifies what would otherwise appear as three separate technical problems.
\end{foldable}

\begin{foldable}
  Treating distributional fidelity as a first-class design objective, rather than a collateral outcome, is the commitment that motivates all three contributions of this thesis.
  Each contribution targets a distinct failure mode, unified by one principle: distillation quality is bounded by signal representativeness, not algorithmic sophistication.
  \S\ref{sec:15-introduction-contributions} describes how each contribution realizes this principle in concrete, experimentally validated form.
\end{foldable}
